\clearpage
\tableofcontents
% \thispagestyle{plain} % 去掉页眉（可选）
\clearpage

\section{Introduction}
The scaling behavior of recommendation models has recently attracted significant attention \citep{OneRec, mtgr, HSTU, Wukong, URM, RankMixer, HLLM, OnePiece}, driven largely by the success of large language models (LLMs) in demonstrating predictable performance gains with increased model size \citep{GPT3, gpt4, deepseek, qwen2.5, mistral-7b, llama3, Llama2, Deepseek-r1}.
However, traditional recommendation models have struggled to exhibit similar scaling laws \citep{SASRec, GRU4Rec, Caser, SGL, survey_seq}.
These systems typically allocate the majority of parameters to large embedding tables for storing user and item representations, while using only inner-product or shallow scoring networks for final predictions.
Such an embedding-heavy design leads to performance plateaus: even as embedding dimensions or table sizes increase, improvements diminish quickly beyond moderate scales \citep{DIN, DCN, DCNv2, DeepFM}.
Generative recommendation offers a fundamental paradigm shift.
By compressing items into sequences of discrete Semantic IDs (SIDs) through quantization techniques \citep{rqkmeans, rqvae}, it uses compact vocabularies and redirects the bulk of parameters toward deep autoregressive Transformers that generate SID sequences \citep{TIGER, LCRec, TOKENRec, OneRec, OnePiece}.
This design enables scaling behaviors more characteristic of language models, where increased depth and capacity translate to consistent performance gains \citep{OneRec}.

Recent industrial deployments have demonstrated the promise of this generative paradigm.
OneRec \citep{OneRec, onerec-tech-v1} achieves significant improvements on Kuaishou's platform with 400 million daily active users, while OneRec-V2 \citep{onerec-tech-v2} further advances through lazy decoder architecture and preference alignment. 
\kxy{OnePiece \citep{OnePiece} integrates LLM-style context engineering and reasoning into both retrieval and ranking models of industrial cascaded pipelines}.
% \wjc{TODO: more prior work from industry.}
However, these results rely on massive proprietary datasets and remain closed-source, leaving critical questions unanswered for the research community:
\begin{itemize}
    \item Do the scaling advantages of generative recommendation transfer to public datasets?
    % \item How do generative models perform across different scales?
    \item What is the minimal post-training recipe needed to achieve strong performance?
\end{itemize}
% Without open implementations and released models, the community cannot regoriouly validate these claims or build upon them systematically.

% In this work, \textbf{we present MiniOneRec, the first fully open-source generative recommendation framework with complete code, training pipelines, and released model checkpoints}.
\kxy{In this work, \textbf{we present MiniOneRec, the first fully open-source generative recommendation framework that provides an end-to-end workflow encompassing SID generation, model Supervised Fine-Tuning (SFT), and recommendation-driven Reinforcement Learning.} The release includes complete source code, reproducible training pipelines, and publicly available model checkpoints.}
We implement SID construction using Residual Quantized Variational Autoencoder (RQ-VAE) \citep{rqvae}, in line with prior work such as TIGER \citep{TIGER, OneRec}.
We post-train Qwen-based (backbone) generative models at multiple scales, ranging from 0.5B to 7B parameters, on public benchmarks from the Amazon Review Dataset.
Following OneRec's successful two-stage post-training pipeline, we adopt SFT on user-item interaction sequences, followed by Group Relative Policy Gradient (GRPO) \citep{Deepseek-r1, deepseekmath} for alignment.
The main contributions of our work are summarized as follows:
\begin{itemize}
    \item \textbf{Validating the Scaling Laws of Generative Recommendation}:
    We present the first systematic investigation of how generative recommendation models scale on public data. Specifically, we train four MiniOneRec variants, spanning 0.5\,B to 7\,B parameters, on the Amazon Review corpus~\citep{amazon} with identical data splits and hyper-parameter settings. Both the final training loss and the held-out evaluation loss consistently decrease as model size grows (see Figure~\ref{fig:scaling};\ref{fig:scaling_eval_loss}), highlighting the superior parameter efficiency of the generative paradigm.
    \item \textbf{Optimizing Post-training Strategies for Generative Recommendation}: 
    We design a lightweight yet comprehensive post-training pipeline that tightly couples full-process SID alignment with reinforced preference optimization. First, we augment the vocabulary with dedicated SID tokens and enforce auxiliary alignment objectives throughout the two-stage optimization process. Second, during the reinforcement learning (RL) phase, we shape the generation process itself: invalid tokens are masked to guarantee that every step produces a legal item, beam search is employed for efficient exploration of diverse candidates, and the reward signal combines rule-based accuracy with a ranking-aware penalty that pushes the model away from hard negatives. The resulting framework, MiniOneRec, offers a concise yet effective recipe for bringing reinforcement learning with value-based ranking into generative recommendation, simultaneously improving candidate diversity and ranking fidelity.
\end{itemize}

The remainder of this report is organized as follows:
Section \ref{bg} reviews background on generative recommendation and the application of LLMs with RL in recommendation systems.
Section \ref{method} presents our proposed MiniOneRec, including task formulation, item tokenization, and post-training pipeline.
Section \ref{training_detail} describes the training details while Section \ref{results} provides comprehensive experimental results.
Section \ref{conclusion} concludes with discussions of future directions.